% Minimal fixture draft for the paperforge smoke tests: two sections, the
% theorem-like zoo, labeled and unlabeled equations, one citation, a
% notation macro defined and used later. Small enough that any failure is
% easy to inspect (review §5-P1).
\documentclass{amsart}
\usepackage{amsmath,amssymb}
\usepackage[colorlinks]{hyperref}
\usepackage[capitalize]{cleveref}

% the one numbering profile the simulator implements: theorem-likes share a
% section-reset counter; equations are numbered globally
\newtheorem{theorem}{Theorem}[section]
\newtheorem{lemma}[theorem]{Lemma}
\newtheorem{definition}[theorem]{Definition}
\newtheorem{remark}[theorem]{Remark}

\newcommand{\wid}{W}
\newcommand{\Zint}{\mathbf{Z}}

\title{A Minimal Fixture Paper}
\author{Alex Author}

\begin{document}

\begin{abstract}
  A tiny paper exercising the paperforge onboarding path: numbering,
  declaration-map candidates, notation, and one citation.
\end{abstract}

\maketitle

\section{Introduction}\label{sec:intro}

We study the widget group $\wid$ over $\Zint$, following \cite{Knuth}.

\begin{definition}\label{def:widget}
  The \emph{widget group} $\wid$ is the trivial group, regarded with the
  discrete topology.
\end{definition}

\begin{theorem}\label{thm:main}
  Every widget is trivial; that is,
  \begin{equation}\label{eq:main}
    \wid = 1.
  \end{equation}
\end{theorem}

\begin{proof}
  Immediate from \cref{def:widget}, since
  \[
    \wid \subseteq 1 .
  \]
\end{proof}

\section{Details}\label{sec:details}

This section proves the supporting lemma used by \cref{thm:main}.

\begin{lemma}\label{lem:helper}
  The widget group $\wid$ has exactly one element.
\end{lemma}

\begin{proof}
  Both inclusions of \cref{eq:main} hold by definition.
\end{proof}

\begin{remark}\label{rem:note}
  Nothing in this paper depends on the choice of $\Zint$.
\end{remark}

\begin{thebibliography}{9}
\bibitem{Knuth} D.~E. Knuth, \emph{The Art of Computer Programming},
  Addison--Wesley, 1968.
\end{thebibliography}

\end{document}
